problem:  
Let \( M^{2n} \) be a compact, simply connected **torus manifold** with an effective \(T^n\)-action and a Riemannian metric of **nonnegative sectional curvature**.  Assume that the orbit space \(Q = M/T^n\) is a simple convex polytope combinatorially equivalent to a **product of simplices**, that is, \(Q \approx \Delta^{m_1} \times \cdots \times \Delta^{m_r}\) with \(m_1 + \cdots + m_r = n\).  Such manifolds are often **quasitoric manifolds**, and many known examples—especially when the cohomology ring is generated in degree two—coincide with **generalized Bott manifolds** arising as successive complex projective bundles.  

Let \(\lambda : \{\text{facets}(Q)\} \to \mathbb{Z}^n\) denote the characteristic function encoding isotropy representations, and let \(H^*(M;\mathbb{Z})\) denote its cohomology ring determined by \((Q, \lambda)\) via the Davis–Januszkiewicz construction.  

> **Proposed Open Problem (Curvature-Rigidity Characterization for Torus Manifolds over Products of Simplices):**  
> Determine whether the condition of **nonnegative sectional curvature** imposes *topological rigidity* on the characteristic data \(\lambda\):  
> - Does nonnegative curvature force \((Q, \lambda)\) to be **equivalent**, up to combinatorial automorphism, to the *trivial characteristic function* corresponding to the product of complex projective spaces?  
> - Equivalently, among all torus manifolds (or quasitoric manifolds) over products of simplices, are the only ones admitting a **\(T^n\)-invariant, nonnegatively curved metric** those that are diffeomorphic to products of rank-one symmetric spaces (e.g. products of \(S^2\) or \(\mathbb{CP}^m\)) or their torus quotients?  
>  
> In a more refined form:  
> *Can nontrivial characteristic twisting (\(\lambda\) non-diagonal) ever coexist with nonnegative sectional curvature?*  
> If not, nonnegative curvature would produce a **geometric rigidity theorem** identifying curvature as a complete topological selector within the class of torus manifolds over products of simplices.

Why is it a "good" problem:  

1. **Conceptually grounded in known frameworks, yet genuinely open.**  
   Existing toric classification results (Davis–Januszkiewicz, Masuda–Panov) completely describe the topology of torus manifolds in combinatorial terms but involve no Riemannian curvature input.  Conversely, geometric rigidity results (Grove–Searle, Wilking) characterize positively curved manifolds by strong symmetry conditions but lack explicit combinatorial control.  
   This problem proposes to join these disjoint frameworks by asking whether curvature enforces *combinatorial triviality* in the characteristic data—an idea unaddressed by existing theory and resistant to known techniques.

2. **Logically precise and technically valid.**  
   The formulation avoids assuming every torus manifold over a product of simplices is a Bott manifold.  Instead, it focuses directly on the characteristic function \(\lambda\) and its geometric constraints, forming a well-defined and credible open question that acknowledges the landscape of partial classification results without contradicting them.

3. **Bridges curvature geometry and toric topology in a deep structural way.**  
   The problem calls for understanding whether curvature—a differential condition on sectional planes—interacts nontrivially with discrete combinatorial data describing isotropy lattices.  Progress would demand integrating tools from Riemannian submersion theory, Alexandrov geometry, rational homotopy constraints (ellipticity), and toric topology, making it a genuine cross-disciplinary frontier.

4. **Extremely simple statement concealing high complexity.**  
   The one-sentence version—  
   “Does nonnegative curvature force a torus manifold over a product of simplices to be a product of symmetric spaces?”—  
   is immediately understandable yet conceals deep nets of algebraic, topological, and geometric relationships.

5. **Potential for a fundamental rigidity principle.**  
   If affirmative, the result would establish a *geometric rigidity for toric topology*: curvature would act as a sieve that removes all nontrivial characteristic twisting.  This would transform the interplay between symmetry rank and topological complexity, possibly linking to broader conjectures on nonnegative curvature and rational ellipticity.

6. **Feasible yet open direction.**  
   No known examples contradict this conjecture; at the same time, no existing geometric argument establishes it.  Its pursuit would illuminate how curvature can constrain torus actions far beyond current symmetry-rank arguments.

In summary, this problem retains conceptual clarity and mathematical precision while probing a **new curvature–combinatorics rigidity phenomenon**: whether *nonnegative sectional curvature annihilates characteristic twisting* in torus manifolds modeled on products of simplices. It thereby epitomizes a “good” problem—simple, profound, plausibly open, and a bridge between two vibrant mathematical worlds.